As a result of libration, studied among others by Johannes Hevelius, more
than half of the Moon's surface can be observed from Earth. Assume that the
observer is geocentric.

\begin{description}
\item[(a)] Estimate $\phi_B$, the maximum angle of libration in latitude. The
  axial tilt (obliquity) of the Moon with respect to its orbital plane is
  $\alpha = 6^\circ 41'$.
\item[(b)] Estimate $\phi_L$, the maximum angle of libration in longitude.
  Assume that the Moon is always aligned with the same side facing towards
  the second focus F2 of its orbit, and that the eccentricity of the Moon's
  orbit $e$ changes between 0.044 and 0.064 on a timescale of several months.
\item[(c)] Estimate the fraction of the Moon's surface which can be seen from
  Earth.
\item[(d)] Calculate how many months (lunations) are needed for an observer
  to see the Moon's surface determined in part (c).
\end{description}
